% native creation stories: base

\subsection{Kabbalah}

Before anything, there was only \textbf{Ein Sof}, the Without End. It had no edge
and no end. It was not a being you could point to. Its infinite light,
\textbf{Or Ein Sof}, filled everything there was, completely, so that there was no
empty place anywhere. There was nothing it could be set against, because there was
nothing else at all.

Ein Sof wished to give. Not because it lacked anything, but because it is the nature
of goodness to pour itself out. To give, there must be someone to receive. So there
had to be a world, and someone in it, standing apart enough to take the gift and, in
time, to give in return.

But a fullness with no edge leaves no room for anything else. A separate thing
cannot exist inside a light that fills all space. So the first act was not an
outpouring. It was \textbf{tzimtzum}, a contraction. Ein Sof withdrew its light away
from a single central point, leaving, for the first time, an empty space, the
\textbf{chalal} or \textbf{tehiru}. Into that cleared room a separate world could
begin. Even in the emptiness a faint trace of the light remained behind, the
\textbf{reshimu}, like a scent in a flask after it is poured out.

Into the empty space Ein Sof sent a single thin ray of light, the \textbf{kav}. The
kav was measured and controlled, unlike the boundless fullness before it. Everything
to come would be built along this one line.

The kav first took the form of \textbf{Adam Kadmon}, the Primordial Man, the pattern
of a whole person laid over the whole of what would be made. From the openings of
this great figure, its eyes and ears and nose and mouth, the light streamed out, and
that streaming became the worlds.

The light organized into the ten \textbf{sefirot}, ten attributes through which the
hidden Ein Sof could act and be known. Chesed and Gevurah, Chokmah and Binah, the
will of Keter and the receiving of Malkhut. These ten repeated across four worlds,
\textbf{Atzilut}, \textbf{Beriah}, \textbf{Yetzirah}, and \textbf{Assiah}, each one
a coarser and more hidden dressing of the light above it. Reality came down like a
chain, link from link, the light growing dimmer and more clothed at every step.

Then came the break. In the world of \textbf{Olam ha-Tohu}, the World of Chaos, the
early vessels, the containers meant to hold each of the ten lights, were strong
enough at the top but too weak below. They stood alone, each grasping the full light
for itself, with no give-and-take between them. Strong light in weak, separate
vessels cannot hold. The lower vessels shattered. This was
\textbf{shevirat ha-kelim}, the breaking of the vessels, the disaster at the root of
the world, the source of all disorder, exile, and death.

When the vessels broke, the shards fell. They became the \textbf{qelippot}, the
shells or husks, the impure \textbf{Sitra Achra}, the Other Side. Clinging to each
falling shard was a holy spark, a \textbf{nitzotz}, trapped now inside the husks.
These husks are the shells that hide the light, the dark side of things. So evil is
not a thing that was made on purpose. It is broken structure with holy light caught
in the wrong place. The light is still holy. It is only imprisoned.

The world we live in is this broken world. It is full of exile and scattered light.

But the work is not finished. Ein Sof planted in the broken world a creature made in
the image of Adam Kadmon, the human soul. By right action and true intention,
\textbf{kavanah}, the soul can find the trapped \textbf{netzotzot} sparks, lift them
out of the husks, and send them home. Each good deed gathers a spark. Each gathered
spark mends a piece of the broken vessels. This is \textbf{tikkun}, the repair.

The story has a clear skeleton. Ein Sof, the boundless source. Tzimtzum, a
drawing-back that makes room. The kav, a single ray of first distinction. A descent
through the four worlds into vessels. Shevirat ha-kelim, the break where the vessels
fail. Holy sparks trapped among the qelippot. A self-aware creature whose labor of
tikkun gathers what was scattered and returns the world to wholeness.

The going-out is creation. The coming-back is history. The one who carries the world
home is the human being, and the whole of that labor is \textbf{tikkun olam}, the
repair of the world.

\subsection{Buddhism}

There is no first moment. There is no maker. There is no one who decided that
anything should exist. Time is \textbf{anadi}, beginningless.

What there is, and has always been, is mind. Not one person's mind, but the bare
fact of awareness, with no beginning. It was never switched on, because there was
never a time before it. In Dzogchen this is the \textbf{gzhi}, the Ground, pure
awareness or \textbf{rigpa}.

This mind has an open and clear nature. Its ultimate nature is \textbf{shunyata},
emptiness, in the sense that it is not a solid thing and has no self-standing
essence. It is bright, in the sense that it knows and lights up whatever appears.
That is what it has always been. Every being holds this awakened potential, called
\textbf{buddha-nature}.

But the mind does not see this about itself. It does not recognize its own open and
clear nature. This single failure to recognize is \textbf{avidya}, ignorance. In
that one failure, everything begins to go wrong, over and over, without ever having
started.

Not seeing what it is, the mind stirs. It reaches. It grasps at what appears, as if
the appearances stood on their own and were separate from it. From that grasping
comes a sense of an inside and an outside, a self in here and a world out there. The
split was never real. Once it is taken as real it feels total.

From that split the appearances pour out. A body to be. A place to be in. Other
beings to meet. A line of time to move along, with a before and an after. The being
itself is only a bundle of processes, the five \textbf{skandhas}, the aggregates of
form, feeling, perception, formations, and consciousness. There is no self,
\textbf{anatman}, to be found among them. None of it is built by a hand. All of it
is the show of a mind that has mistaken its own light for a world standing outside
it.

The grasping leaves marks. \textbf{Karma}, intentional action, plants seeds in the
mindstream, the \textbf{alaya} or storehouse consciousness, and each seed ripens
later into the next situation, the next body, the next place. This turning is laid
out as the twelve \textbf{nidanas}, the links of dependent arising,
\textbf{pratityasamutpada}, ignorance to formations to consciousness and on around
the wheel. So the appearances do not rise once. They keep rising, shaped by what
came before, life after life, with no finish line and no opening line. This endless
round is \textbf{samsara}.

This is the whole of it. A self-built dream. No dreamer apart from the dreaming, no
maker apart from the mind, no first night when the dream began.

Nothing decides that the worlds should be. There is no creator god and no first
cause. The worlds are an effect with no author, only a long chain of causes that
never starts. The turning of forming and dissolving runs in \textbf{kalpas}, immense
cycles of formation, duration, destruction, and emptiness, going round and round,
and any point you pick already has a point before it. The starting condition is not
stuff and it is not a maker. It is awareness that fails to recognize its own nature.
The world comes from grasping, not from a word spoken over an empty void.

And the same mind that spins the dream can wake from it. That waking is
\textbf{bodhi}, and a mind that has fully woken is a \textbf{buddha}, one who has
cleared the obscurations and recognized its own nature. The waking changes nothing
about the fact that there was no first night. The world is real enough to suffer in
and to wake up from. It is just not made, and it never began.
